\section{Test Setup and Requirements}

\subsection{Evaluation Criteria}
\begin{table}[H]
    \centering
    \small
    \begin{tabular}{l c p{7.2cm}}
        \toprule
        \textbf{Requirement} & \textbf{Target} & \textbf{Rationale} \\
        \midrule
        Input Voltage Range & \SI{9}{\volt} to \SI{16}{\volt} & Representative operating envelope for the subsystem. \\
        Output Accuracy & $\pm$5\% & Ensures acceptable regulation under nominal load. \\
        Peak Temperature & $< \SI{90}{\celsius}$ & Maintains component margin and reliability. \\
        Startup Behavior & No fault or latch-up & Prevents unsafe transients during bring-up. \\
        \bottomrule
    \end{tabular}
    \caption{Representative evaluation targets}
    \label{tab:requirements}
\end{table}

\subsection{Bench Configuration}
\begin{itemize}
    \item programmable power supply for input stimulus;
    \item electronic load or representative passive load;
    \item oscilloscope for transient capture;
    \item multimeter or power analyzer for steady-state measurements;
    \item thermal camera or thermocouple instrumentation when temperature matters.
\end{itemize}

\warningbox{
Before copying measured data into the report, confirm units, channel names, and
test conditions. Most review confusion comes from screenshots that are not tied
to a clearly stated setup.
}

\subsection{Bill of Materials Snapshot}
\makecomponenttable{Representative Critical Components}{
    U1 & Main controller / regulator under evaluation & ABC1234 & Generic Semiconductor \\
    Q1 & Switching device for the primary load path & QSW-45A & Generic Power Devices \\
    D1 & Protection or flyback diode & DFLY-40V & Generic Rectifier Co. \\
}
